\chapter[综合实战案例：智能制造知识管理系统]
  {综合实战案例：智能制造\\知识管理系统}
\chapterart{ch09}{从双栈样例到 Semantica 单一执行语义}

\begin{introbox}
\textbf{【导读】}
最后一章把前八章的概念、CQ、语言、查询、规则和治理意识汇入一个制造业综合案例。
旧版的学习路线是分别跑通 Java/Jena 查询服务与 Python/owlready2/Pellet 推理器；
本版保留这些代码片段作为书史和技术比较，却不再把它们当成运行入口。
现在的学习路线是「查看 manifest—运行具名 scenario—核对 exact oracle—单独验证 release」：
所有可执行语义都由 Semantica 内建包承载，未实现的 Manchester、SWRL、DL 与 CQ
明确阻断，不能靠另一套本地工具链偷偷变绿。
\end{introbox}

\chapterbinding{semantica.chapter\_packages.vol1.ch09}{partial}{blocked}
{章合同、CQ1--CQ7、SPARQL 集合、Manchester 来源本体、结构化 SWRL 规则、
CQ1 正例与单因反例均已迁入；\texttt{OE-V1-CH09-SCN-CQ01-001} 由原生 runner 执行。}
{CQ2--CQ7 缺少书内 ABox exact oracle；本章没有 SHACL；当前 runtime 不解析
Manchester 或一般 SWRL，也不执行完整 DL/tableau。旧 Jena/owlready2 代码不属于运行时。}

\section{系统全景：书是规格，包是唯一实现}

这套系统不再按编程语言分成两套权威实现，而按责任分成五层：

\begin{enumerate}
\item \textbf{可读规格层}：本章叙事、CQ、设计理由、练习与历史实现比较；
\item \textbf{包身份层}：manifest、contract、asset hash、scenario registry 与 capability gap；
\item \textbf{语义资产层}：本体来源、查询、规则、夹具与 oracle，均在 ch09 包内；
\item \textbf{执行层}：\texttt{SemanticPackageRunner} 委托同一个 \texttt{SemanticRuntime}；
\item \textbf{证据与发布层}：输入/输出/运行时哈希、PROV、receipt 与独立 release verdict。
\end{enumerate}

这种划分把三个常被混在一起的权威拆开：包定义的是\textbf{语义权威}；
车间台账、测量与受控记录提供\textbf{事实权威}；具名负责人和有权限的系统承担
\textbf{决策与动作权威}。查询返回一台设备，绝不等于系统获准启动它。

包内 \texttt{manufacturing} 资产保留了 Manchester Syntax 的类、公理与注释。
下列节选仍是理解设计意图的重要“石头”，但当前 runtime 不解析 Manchester，
因此它是 \texttt{ontology\_source}，不是已执行的 DL 证明：

\input{fragments/snip-ch09-axioms}
\assetsource{ch09}{manufacturing}

\section{CQ1 的可复算旅程}

CQ1 问“哪些设备可以加工指定材料”。在 Semantica 中，一条可复算主张必须同时固定：

\begin{itemize}
\item 查询资产 \texttt{cq01}；
\item 正例 \texttt{cq01-positive} 与只改变一个条件的
      \texttt{cq01-single-fault-negative}；
\item 场景 \texttt{OE-V1-CH09-SCN-CQ01-001}；
\item exact multiset oracle：正例返回“数控车床一号”“立式铣床二号”，反例返回 0 行；
\item package、asset、输入、输出与 runtime 身份绑定的执行结果。
\end{itemize}

旧版 \texttt{QueryService.getEquipmentForMaterial} 的替代入口不是另一段书内代码，
而是同一 runner 中的 \texttt{execute\_chapter\_query}。这保证 CLI、Python 与 MCP
只是适配面，不会各自解释一套 SPARQL 语义。

主运行路径是 ontology-engineering 的统一 semantic engagement 入口。它自动核验并注入
正式 source lock 中的 runtime commit、版本与 wheel SHA-256；binding 与 task fixture
必须按 \texttt{references/semantic-engagement-contract.md} 建立，不能手抄运行时身份。

\begin{codebox}
\codeline{runtime/.venv/bin/python\ scripts/semantic\_engagement.py\ discover}
\codeline{runtime/.venv/bin/python\ scripts/semantic\_engagement.py\ run\ \textbackslash{}}
\codeline{\ \ --binding\ /path/to/package-binding.json\ --task\ /path/to/task-envelope.json\ \textbackslash{}}
\codeline{\ \ --scenario\ OE-V1-CH09-SCN-CQ01-001}
\codeline{runtime/.venv/bin/python\ scripts/semantic\_engagement.py\ open\ \textbackslash{}}
\codeline{\ \ --binding\ /path/to/workspace-binding.json\ --task\ /path/to/task-envelope.json\ \textbackslash{}}
\codeline{\ \ --workspace\ /path/to/semantica-managed-registry}
\end{codebox}

\texttt{run} 回答“声明场景是否符合 oracle”；\texttt{verify} 还检查来源、能力、
内容与发布证据。当前即使 CQ1 复算成功，release 仍应是 \texttt{blocked}。
\texttt{propose/commit/verify/history/promote} 继续消费合同定义的 delta、candidate 与
gate-evidence fixtures；参数以统一入口各子命令的 \texttt{--help} 为准。
原生 \texttt{semantica package show/run/verify} 只保留为底层诊断接口，不是主运行路径。

\section{为什么 CQ2--CQ7 不能跟着 CQ1 一起变绿}

包内已经有七条 SPARQL 查询，并不等于七条 CQ 都通过。CQ2--CQ7 缺少与书中问题
一一对应的 ABox 正例、单因反例与精确预期绑定；仅证明“查询可解析”或“返回了若干行”
都不足以构成验收。后续每补一条 CQ，都应独立给出：

\begin{enumerate}
\item 最小输入图与适用范围；
\item 精确变量、行数、多重集和排序是否属于 oracle；
\item 单因反例，以及 Unknown 与 false 的区分；
\item 前一发布版回归；
\item scenario receipt 与仍然独立的 package release verdict。
\end{enumerate}

同理，五条 SWRL 规则已经结构化迁入 \texttt{engineering-rules}，但当前
\texttt{SemanticRuntime} 不解析或执行一般 SWRL 与 built-ins。不能把规则文本存在、
人工看起来合理或第三方引擎曾经跑过，写成 Semantica 已证明的结论。

\section{旧双栈为什么退出正文执行链}

旧版曾用 Java/Jena 演示查询服务，用 Python/owlready2/Pellet 演示推理工作流。
它们帮助过读者认识“装载—查询—推理”的传统分层，却也带来两套 parser、profile、
依赖版本、错误处理和结果证据。若继续排印安装与运行片段，读者很容易在 Semantica
报告 unsupported 时绕去另一套引擎，最终得到无法与 package identity 和 receipt
绑定的平行答案。

因此本版删除旧 Java/Python 代码框、Maven/pip 安装步骤与本地运行说明，只在这里
记录架构演化事实。需要比较生态时可阅读第 4、5 章的原理性介绍；需要复算本章主张时
只有 ch09 package/runner 一个入口。书史是“石头”，旧运行时不是。

\section{从教学案例走向工程发布}

将本章扩成现实系统时，建议按以下顺序推进：

\begin{enumerate}
\item 为 CQ2--CQ7 逐条建立最小正例、单因反例与 exact oracle；
\item 增加本章自有 SHACL，显式区分开放世界推理与封闭式完整性检查；
\item 若确需 Manchester、OWL DL 或 SWRL，先冻结支持 profile 与合规测试，再实现；
\item 对每次语义变更执行 snapshot/diff、旧 CQ 回归、PROV 和内容绑定 receipt；
\item 将真实 ABox 接入、写操作、设备动作和补偿交给具名且有权限的执行面；
\item 只有所有必需证据齐备并通过独立发布门禁时，才改变 release status。
\end{enumerate}

\begin{noteblock}
\textbf{综合案例的完成标准不是“程序启动了”。}
完成意味着：书中问题能定位到唯一 package 与 scenario；输入、输出和 runtime 可复核；
不支持的能力被阻断；事实与决定来自有权来源；旧版本回归不被静默破坏。
当前 ch09 仍是 partial/blocked，这个红灯本身就是可信工程的一部分。
\end{noteblock}

\section{综合案例也是行业本体的炼化入口}

现实项目不应另建一套“从经验长本体”的旁路。每次使用本章方法时，项目用严格
task envelope 交代目的、能力需求和带哈希证据，用 project binding 锁定 Semantica
workspace、baseline、事实权威、决定权威和允许推进的生命周期。工作结束必须同时看
三份账：工程问题是否解决，现有 package 的执行/回归/receipt/release 是否闭合，
以及这次实践是否真的带来了跨项目可复用的 delta。

若答案是 candidate，完整变更包要覆盖 ontology、CQ、SHACL、query、rule、四类案例、
contract、provenance 与独立的 book impact；然后按 proposed、committed、
regression passed、release complete、promoted 顺序推进。任何阶段的 exit code、一次
CQ 通过或一个“看起来合理”的解释都不能跳级。项目事实仍留在项目证据层，只有经过
作用域判断、冲突裁决、旧知识回归与有权晋升的结构才进入行业 registry。

\section{本章资产指引}

\begin{center}
\begin{tabularx}{\linewidth}{@{}lXX@{}}
\toprule
\textbf{包内资产/场景} & \textbf{用途} & \textbf{当前边界} \\
\midrule
\texttt{manufacturing} & Manchester 来源本体与设计说明 & 保留，不解析、不声称 DL 通过 \\
\texttt{engineering-rules} & 从来源本体结构化提取的五条 SWRL & 保留，不执行一般 SWRL/built-ins \\
\texttt{cq01}--\texttt{cq07} & 制造领域查询正本 & 仅 CQ1 绑定 exact oracle \\
\texttt{cq01-positive} / 单因反例 & CQ1 受控输入 & 原生 scenario 使用 \\
\texttt{OE-V1-CH09-SCN-CQ01-001} & CQ1 精确多重集验收 & 可运行；不等于 package release \\
旧 Java/Python 片段 & 实现史与技术比较 & 不在 runtime，不是 receipt \\
\bottomrule
\end{tabularx}
\end{center}

到这里，全书形成了统一闭环：第 1 章定义目的与边界，第 2--5 章解释语义与方法，
第 6--8 章展示应用、数据与智能体控制，第 9 章用一个诚实保留红灯的 Semantica 包
把它们汇合。下一次扩展从包内新增可判定证据开始，而不是从书旁复制一份运行时开始。
